%% ma: L12-B07-0006
%% lop: 12
%% chuong: 2
%% bai: 7
%% dang: 12.7.5
%% loai: TN4
%% muc: NB
%% dap_an: "A"
%% nguon: "LaTeX giáo viên gửi 21/10/2026 (M05, Vectơ và hệ trục tọa độ trong không gian, 107 câu), Phần I Câu 1"
%% kien_thuc: "Bài 7 (tọa độ của vectơ)"
%% trang_thai: cho-duyet
%% ly_do: "Dạng mới đề xuất 12.7.5"
\begin{ex}
Trong không gian $Oxyz$, cho vectơ $\vec{a}=2\vec{i}-3\vec{j}+5\vec{k}$. Tọa độ của vectơ $\vec{a}$ là
\choice{\True $(2;-3;5)$}{$(2;5;-3)$}{$(-3;2;5)$}{$(2;0;5)$}
\loigiai{$\vec{a}=x\vec{i}+y\vec{j}+z\vec{k}\Leftrightarrow\vec{a}=(x;y;z)$ nên $\vec{a}=(2;-3;5)$.}
\end{ex}
